% Part: propositional-logic
% Chapter: propositional-logic
% Section: preliminaries

\documentclass[../../../include/open-logic-section]{subfiles}

\begin{document}

\olfileid{pl}{syn}{pre}

\olsection{ప్రాథమిక ఫలితాలు}

\begin{thm}[\emph{\tetoken{సూత్రాలపై}{!!{formula}s} ఆగమన సూత్రం}]
  \ollabel{thm:induction}
  ఏదైనా ధర్మం~$P$ అన్ని పరమాణు \tetoken{సూత్రాలకు}{!!{formula}s} ఉండి, అది కింది
  విధంగా ఉందనుకుందాం:
  \begin{enumerate}
    \tagitem{prvNot}{$!A$కు అది ఉన్నప్పుడల్లా $\lnot !A$కు కూడా
      ఉంటుంది;}{}
    \tagitem{prvAnd}{$!A$, $!B$లకు అది ఉన్నప్పుడల్లా
      $(!A \land !B)$కు ఉంటుంది;}{}
    \tagitem{prvOr}{$!A$, $!B$లకు అది ఉన్నప్పుడల్లా
      $(!A \lor !B)$కు ఉంటుంది;}{}
    \tagitem{prvIf}{$!A$, $!B$లకు అది ఉన్నప్పుడల్లా
      $(!A \lif !B)$కు ఉంటుంది;}{}
    \tagitem{prvIff}{$!A$, $!B$లకు అది ఉన్నప్పుడల్లా
      $(!A \liff !B)$కు ఉంటుంది;}{}
  \end{enumerate}
  అప్పుడు $P$ అన్ని \tetoken{సూత్రాలకు}{!!{formula}s} ఉంటుంది.
\end{thm}

\begin{proof}
  $P$ ధర్మం గల అన్ని \tetoken{సూత్రాలు}{!!{formula}s} సంగ్రహాన్ని $S$ అనుకుందాం.
  స్పష్టంగా $S \subseteq \Frm[L_0]$. $S$,
  \olref[fml]{defn:formulas}లోని అన్ని షరతులను సంతృప్తిపరుస్తుంది:
  అది అన్ని పరమాణు \tetoken{సూత్రాలను}{!!{formula}s} కలిగి ఉండి, \tetoken{తార్కిక సంయోజకాలు}{!!{operator}s} కింద
  సంవృతంగా ఉంటుంది. అటువంటి అతి చిన్న వర్గం $\Frm[L_0]$ కనుక
  $\Frm[L_0] \subseteq S$. కాబట్టి $\Frm[L_0] = S$; ప్రతి సూత్రానికీ
  $P$ ధర్మం ఉంటుంది.
\end{proof}

\begin{prop}\ollabel{prop:balanced}
   $\Frm[L_0]$లోని ఏ \tetoken{సూత్రం}{!!{formula}} అయినా \emph{సమతుల్యం}; అంటే దానిలో
   ఎన్ని ఎడమ కుండలీకరణలు ఉంటాయో అన్నే కుడి కుండలీకరణలు ఉంటాయి.
\end{prop}

\begin{prob}
\olref[pl][syn][pre]{prop:balanced}ను నిరూపించండి.
\end{prob}

\begin{prop} \ollabel{prop:noinit}
\tetoken{ఒక సూత్రం}{!!a{formula}} యొక్క ఏ నిజ ఆద్య ఖండమూ \tetoken{ఒక సూత్రం}{!!a{formula}} కాదు.
\end{prop}

\begin{prob}
  \olref[pl][syn][pre]{prop:noinit}ను నిరూపించండి.
\end{prob}

\begin{prop}[ఏకైక పఠనీయత]
$ \Frm[L_0]$లోని ఏ \tetoken{సూత్రం}{!!{formula}}~$!A$కైనా కింది రూపాల్లో కచ్చితంగా
ఒకే ఒక వాక్యనిర్మాణ విశ్లేషణ ఉంటుంది:
\begin{enumerate}
\tagitem{prvFalse}{$\lfalse$.}{}

\tagitem{prvTrue}{$\ltrue$.}{}

\item ఏదో $\Obj p_n \in  \PVar$కు $\Obj p_n$.

\tagitem{prvNot}{ఏదో \tetoken{సూత్రం}{!!{formula}}~$!B$కు $\lnot !B$.}{}

\tagitem{prvAnd}{ఏవో \tetoken{సూత్రాలు}{!!{formula}s} $!B$, $!C$లకు
$(!B \land !C)$.}{}

\tagitem{prvOr}{ఏవో \tetoken{సూత్రాలు}{!!{formula}s} $!B$, $!C$లకు
$(!B \lor !C)$.}{}

\tagitem{prvIf}{ఏవో \tetoken{సూత్రాలు}{!!{formula}s} $!B$, $!C$లకు
$(!B \lif !C)$.}{}

\tagitem{prvIff}{ఏవో \tetoken{సూత్రాలు}{!!{formula}s} $!B$, $!C$లకు
$(!B \liff !C)$.}{}
\end{enumerate}
అంతేకాక ఈ వాక్యనిర్మాణ విశ్లేషణ \emph{ఏకైకం}.
\end{prop}

\begin{proof}
$!A$పై ఆగమనం ద్వారా నిరూపిస్తాం. ఉదాహరణకు, $!A$కు
$(!B \lif !C)$, $(!B' \lif !C')$ అనే రెండు భిన్న పఠనాలు ఉన్నాయని
అనుకుందాం. అప్పుడు $!B$, $!B'$ ఒకటే కావాలి (లేకపోతే ఒకటి మరొకదానికి
నిజ ఆద్య ఖండమవుతుంది). కాబట్టి $!A$ యొక్క రెండు పఠనాలు భిన్నమైతే,
ఒకే సంకేత క్రమానికి $!C$, $!C'$లు భిన్న విశ్లేషణలు కావడమే కారణం; కానీ
ఆగమన పరికల్పన వల్ల అది అసాధ్యం.
\end{proof}


\begin{defn}[ఏకరీతి ప్రతిస్థాపన]
$!A$, $!B$లు \tetoken{సూత్రాలు}{!!{formula}s}, $\Obj p_i$ ఒక \tetoken{ప్రతిజ్ఞావాక్య చరం}{!!{propositional
variable}} అయితే, $!A$లోని $\Obj p_i$ ప్రతి సంభవాన్ని $!B$ సంభవంతో
మార్చడం వల్ల వచ్చే ఫలితాన్ని $\Subst{!A}{!B}{\Obj p_i}$ సూచిస్తుంది.
అలాగే $\Obj p_1$, \dots,~$\Obj p_n$లను \tetoken{సూత్రాలు}{!!{formula}s} $!B_1$,
\dots,~$!B_n$లతో ఏకకాలంలో ప్రతిస్థాపించడాన్ని
$\SSubst{!A}{\subst{!B_1}{\Obj p_1},\dots,\subst{!B_n}{\Obj p_n}}$తో
సూచిస్తాం.
\end{defn}

\begin{prob} కింద ఉన్న ఐదు \tetoken{సూత్రాలలో}{!!{formula}s} ప్రతి ఒక్కదానికి, ఆ
 \tetoken{సూత్రాన్ని}{!!{formula}} \( \Subst{!A}{!B}{\Obj p_i} \) అనే ప్రతిస్థాపనగా
 రాయవచ్చో లేదో నిర్ణయించండి; ఇక్కడ \( !A \) అనేది (i)
 \( \Obj p_0 \); (ii) \( ( \lnot \Obj p_0 \land \Obj
 p_1) \); మరియు (iii) \( ( ( \lnot \Obj p_0 \lif \Obj p_1 ) \land \Obj
 p_2 ) \). ప్రతి సందర్భంలో సంబంధిత ప్రతిస్థాపనను పేర్కొనండి.
  \begin{enumerate}
    \item \( \Obj p_1 \)
    \item \( ( \lnot \Obj p_0 \land \Obj p_0 ) \)
    \item \( ( ( \Obj p_0 \lor \Obj p_1 ) \land \Obj p_2 )  \)
    \item \( \lnot ( ( \Obj p_0 \lif \Obj p_1 ) \land \Obj p_2 ) \)
    \item \( (( \lnot ( \Obj p_0 \lif \Obj p_1 ) \lif ( \Obj p_0 \lor \Obj p_1 )) \land \lnot ( \Obj p_0 \land \Obj p_1 )) \)
  \end{enumerate}
\end{prob}

\begin{prob}
ఆగమనం ద్వారా $\Subst{!A}{!B}{p}$కు గణితపరంగా కచ్చితమైన నిర్వచనం
ఇవ్వండి.
\end{prob}

\end{document}
